\documentclass{article}
\usepackage[T1]{fontenc}
\usepackage{lmodern}
\usepackage{graphicx}
\usepackage{amsmath,amssymb}
\usepackage[a4paper,margin=2cm]{geometry}
\usepackage[absolute,overlay]{textpos}
\setlength{\TPHorizModule}{1mm}
\setlength{\TPVertModule}{1mm}
\usepackage[indonesian]{babel}
\usepackage{hyperref}
\setlength{\emergencystretch}{2em}
% unit: Halaman 50

\begin{document}

% === Halaman 50 (python/native) ===

50

def dekripsi\_aes(cipherteks, kunci):

    kunci = kunci.encode()                     \# Kodekan kunci dari string menjadi byte

    cipherteks = base64.b64decode(cipherteks)  \# Dekodekan kembali cipherteks dari base64 ke kode byte

    iv = cipherteks[:16]                      \# Ambil iv pada bagian awal cipherteks sepankang 16 byte

    AEScipher = AES.new(kunci, AES.MODE\_CBC, iv)        \# Defenisikan AES mode CBC 

    plainteks = AEScipher.decrypt(cipherteks[BLOCK\_SIZE:])           \# Lakukan dekripsi cipherteks

    plainteks = unpad(plainteks, BLOCK\_SIZE)            \# Buang padding

    return plainteks  

def encrypt\_file(input\_file, output\_file, key):

    with open(input\_file, 'rb') as f:

         plaintext = f.read()

    ciphertext = enkripsi\_aes(plaintext, key)

    with open(output\_file, 'wb') as f:

         f.write(ciphertext)

def decrypt\_file(input\_file, output\_file, key):

    with open(input\_file, 'rb') as f:

         ciphertext = f.read()

    plaintext = dekripsi\_aes(ciphertext, key)

    with open(output\_file, 'wb') as f:

        f.write(plaintext)

\end{document}
